\subsection{Gelfand 变换}

设 A 是一个交换 Banach 代数。回顾：对每个 \(x \in A\)，我们用 \(\mathcal{G}_{\mathrm{A}}(x)\) 或 \(\mathcal{G}(x)\) 表示函数 \(\chi \mapsto \chi(x)\)（定义在 \(\mathsf{X}(\mathsf{A})\) 上）；\(\mathcal{G}(x)\) 称为 x 的 Gelfand 变换，而映射 \(x \mapsto \mathcal{G}(x)\) 称为 Gelfand 变换（参见 I, p.~7, 定义 5）。于是按定义有：

\[
\mathcal {G} (x) (\chi) = \chi (x).
\]

例. --- 1) 设 X 是一个局部紧拓扑空间，并考虑交换 Banach 代数 \(\mathcal{C}_{0}(\mathrm{X})\) (I, p.~17, 例 3 及 I, p.~31, 第 2 小节)。由 I, p.~32, 推论 1，特征标空间 \(\mathsf{X}(\mathcal{C}_{0}(\mathbf{X}))\) 借助于把元素 \(x \in \mathbf{X}\) 联系到特征标 \(f \mapsto f(x)\)（\(\mathcal{C}_{0}(\mathbf{X})\) 的特征标）的映射而与 X 视为同一，于是 \(\mathcal{C}_{0}(\mathbf{X})\) 的 Gelfand 变换与恒等映射视为同一。

\begin{enumerate}
\def\labelenumi{\arabic{enumi})}
\setcounter{enumi}{1}
\tightlist
\item
  设 \(n \geqslant 0\) 是一个整数。设 \(A_n\) 是由函数 \(f: [0,1] \to K\)（在 \([0,1]\) 上具有直到 \(n\) 阶的连续导数）所成的代数。赋以范数
\end{enumerate}

\[
\| f \| = \sum_ {k = 0} ^ {n} \frac {1}{k !} \sup _ {0 \leqslant t \leqslant 1} | f ^ {(k)} (t) |,
\]

它是一个 Banach 代数 (I, p.~18, 例 4)。对每个 n，特征标空间 \(\mathsf{X}(\mathsf{A}_{n})\) 与 {[}0,1{]} 视为同一，而 G 与 \(A_{n}\) 到 \(\mathcal{C}([0,1])\) 中的包含映射视为同一（参见 I, p.~144, 例 1）。

\begin{enumerate}
\def\labelenumi{\arabic{enumi})}
\setcounter{enumi}{2}
\item
  设 \(\Delta\) 是由复数 \(z\)（满足 \(|z| \leqslant 1\)）组成的圆盘，并设 A 是由 \(\Delta\) 上连续且在 \(\Delta\) 的内部解析的函数所成的复 Banach 代数，赋以范数 \(\|f\| = \sup_{z \in \Delta} |f(z)|\) (I, p.~20, 例 9)。则 \(\mathsf{X}(\mathsf{A})\) 与 \(\Delta\) 视为同一，而 \(\mathcal{G}\) 与 A 到 \(\mathcal{C}(\Delta)\) 中的包含映射视为同一（参见 I, p.~193, 习题 6）。
\item
  考虑绝对收敛 Fourier 级数所成的复 Banach 代数 A (I, p.~19, 例 8)。对单位圆 U 的每个元素 \(u\)，映射 \(f \mapsto f(u)\) 是 A 的一个特征标 \(ev_{u}\)。若 \(f_{0} \in A\) 是 U 的恒等映射，则有 \(ev_{u}(f_{0}) = u\)，因而映射 ev: \(u \mapsto ev_{u}\)（从 U 到 X(A)）是单射；它是连续的。设 \(\chi \in \mathsf{X}(\mathsf{A})\)。我们有 \(\|f_{0}\| = \|f_{0}^{-1}\| = 1\)，因而 \(|\chi(f_{0})| \leqslant 1\) 且 \(|\chi(f_{0})^{-1}| \leqslant 1\)。这表明 \(\chi(f_{0}) \in \mathbf{U}\)，并且存在 \(u \in U\) 使得 \(\chi(f_{0}) = \mathrm{ev}_{u}(f_{0})\)。由于 \(\{f_{0}, f_{0}^{-1}\}\) 拓扑生成幺代数 A，有 \(\chi = ev_{u}\)。于是，映射 ev 是从 U 到 X(A) 上的一个同胚，借此把这两个空间视为同一。A 的 Gelfand 变换于是与 A 到 \(\mathcal{C}(\mathbf{U})\) 中的包含映射视为同一。
\end{enumerate}

由于 A 同构于 Banach 代数 \(\mathrm{L}^{1}(\mathbf{Z})\) (I, p.~19, 例 8)，空间 \(\mathsf{X}(\mathrm{L}^{1}(\mathbf{Z}))\) 与 U 视为同一，并且对每个元素 \((c_{n})\in\mathrm{L}^{1}(\mathbf{Z})\)，Gelfand 变换 \(\mathcal{G}_{\mathrm{L}^{1}(\mathbf{Z})}((c_{n}))\) 与 U 上的函数 \(u\mapsto\sum_{n\in\mathbf{Z}}c_{n}u^{n}\) 视为同一。

\begin{enumerate}
\def\labelenumi{\arabic{enumi})}
\setcounter{enumi}{4}
\tightlist
\item
  设 \(\Delta\) 是由复数 \(z\)（满足 \(|z| \leqslant 1\)）组成的单位圆盘。它在 \(\mathbf{C}\) 中的边界是 \(\mathbf{U}\)。设 A 是由复函数 \(f\)（\(\mathbf{U}\) 上的函数，存在连续函数 \(\widetilde{f} \in \mathcal{C}(\Delta)\) 延拓 \(f\) 且在 \(\mathring{\Delta}\) 中解析）所成的 Banach 代数，赋以范数 \(\|f\| = \sup_{z \in \mathbf{U}} |f(z)|\)。于是由最大模原理 (VAR, R1, p.~30, 3.3.7)，有 \(\|f\| = \sup_{z \in \Delta} |\widetilde{f}(z)|\)，因而 A 与 I, p.~20, 例 9 中的代数重合。集合 \(\mathsf{X}(\mathsf{A})\) 与 \(\Delta\) 视为同一，且若 \(f \in \mathsf{A}\)，则映射 \(\mathcal{G}(f)\) 与 \(f\) 到 \(\Delta\) 上的、在 \(\mathring{\Delta}\) 中解析的连续延拓视为同一（参见 I, p.~193, 习题 6）。
\end{enumerate}

命题 5. --- 设 A 是一个交换 Banach 代数。对每个 \(x \in \mathbf{A}\)，函数 \(\mathcal{G}(x)\) 属于交换 Banach 代数 \(\mathcal{C}_0(\mathsf{X}(\mathbf{A}))\)（由 \(\mathsf{X}(\mathbf{A})\) 上在无穷远处为 0 的连续函数所成）。

按定义（参见 I, p.~9, 第 7 小节），函数 \(\mathcal{G}_{\mathrm{A}}^{\prime}(x)\colon \chi \mapsto \chi (x)\) 在 \(\mathsf{X}'(\mathsf{A})\) 上连续且在 0 处取值 0。由于 \(\mathsf{X}'(\mathsf{A})\) 依 I, p.~29, 推论 1 与 \(\mathsf{X}(\mathsf{A})\) 的 Alexandroff 紧化视为同一，命题得证。

命题 6. --- 设 A 是一个交换 Banach 代数，并设 \(x \in A\)。

\begin{enumerate}
\def\labelenumi{\alph{enumi})}
\item
  \(\mathcal{G}(x)\) 的值集与 \(\{0\}\) 的并等于 \(\mathrm{Sp}_{\mathrm{A}}^{\prime}(x)\)；
\item
  若 A 有单位元，则 \(\mathcal{G}(x)\) 的值集就是 \(\mathrm{Sp}_{\mathrm{A}}(x)\)。特别地，要使 \(x\) 可逆，必须且只需 \(\mathcal{G}(x)\) 不取值 0。
\end{enumerate}

假定 A 有单位元。我们知道，对每个 \(\chi \in \mathsf{X}(\mathsf{A})\)，有 \(\chi(x) \in \mathrm{Sp}_{\mathsf{A}}(x)\)。反之，设 \(\lambda \in \mathrm{Sp}_{\mathsf{A}}(x)\)。则 \(x - \lambda\) 不可逆，因而属于 A 的某个极大理想。于是存在 \(\chi \in \mathsf{X}(\mathsf{A})\) 使得 \(\chi(x - \lambda) = 0\) (I, p.~30, 定理 2)，由此得 \(b\)）。

过渡到一般情形。设 \(\widetilde{\mathsf{A}}\) 是由 A 添加单位元所得到的 Banach 代数；它是交换的。集合 \(\mathrm{Sp}_{\widetilde{\mathsf{A}}}^{\prime}(x)\) 等于 \(\mathrm{Sp}_{\widetilde{\mathsf{A}}}^{\sim}(x)\)，即 \(\mathcal{G}_{\widetilde{\mathsf{A}}}^{\sim}(x)\) 在 \(\mathsf{X}(\widetilde{\mathsf{A}}) = \mathsf{X}'(\mathsf{A})\) 上的值集。由此得 \(a\)）。

例. --- 考虑绝对收敛 Fourier 级数所成的 Banach 代数 A (例 4)。命题 6, b) 蕴涵：若 \(\varphi\) 是单位圆 U 上一个具有绝对收敛 Fourier 级数的函数，且 \(\varphi\) 不取值 0，则函数 \(1/\varphi\) 也具有绝对收敛 Fourier 级数（“Wiener 定理”）。

命题 7. --- 设 A 是一个交换 Banach 代数。

\begin{enumerate}
\def\labelenumi{\alph{enumi})}
\item
  Gelfand 变换 \(\mathcal{G}\) 定义一个从 A 到 \(\mathcal{C}_0(\mathsf{X}(\mathsf{A}))\) 的态射，使得 \(\| \mathcal{G}(x)\| = \varrho (x)\leqslant \| x\|\) 对一切 \(x\in \mathrm{A}\) 成立；
\item
  要使 Gelfand 变换 \(\mathcal{G}\) 是等距的，必须且只需 \(\|x^{2}\|=\|x\|^{2}\) 对一切 \(x\in\mathrm{A}\) 成立。
\end{enumerate}

依 I, p.~9, 第 7 小节及 I, p.~37, 命题 5，映射 \(\mathcal{G}\) 是从 A 到 \(\mathcal{C}_0(\mathsf{X}(\mathsf{A}))\) 的态射；依命题 6 及 I, p.~28, 推论 5，它满足 \(\| \mathcal{G}(x)\| = \varrho (x)\)。断言 \(b)\) 由 \(a)\) 与 I, p.~21, 注 1 得出。

推论. --- 设 A 是一个 Banach 代数，并设 x 与 y 是 A 中可交换的元素。

\begin{enumerate}
\def\labelenumi{\alph{enumi})}
\item
  我们有 \(\varrho(xy) \leqslant \varrho(x)\varrho(y)\) 及 \(\varrho(x + y) \leqslant \varrho(x) + \varrho(y)\)；
\item
  若 y 是拟幂零元，则 \(\mathrm{Sp}_{\mathrm{A}}^{\prime}(x)=\mathrm{Sp}_{\mathrm{A}}^{\prime}(x+y)\)；此外若 A 是幺的，则 \(\mathrm{Sp}_{\mathrm{A}}(x)=\mathrm{Sp}_{\mathrm{A}}(x+y)\)。
\end{enumerate}

考虑由 A 添加单位元所得到的 Banach 代数，首先化归到代数 A 是幺的情形。然后，考虑由 x 与 y 生成的 A 的闭全子代数，就化归到 A 是交换且幺的情形。这时断言 (a) 是命题 7(a) 的推论，而断言 (b) 由 I, p.~37, 命题 6 及 I, p.~28, 推论 5 得出。

命题 8. --- 设 A 是一个交换 Banach 代数。下列四个集合相等：

\begin{enumerate}
\def\labelenumi{(\roman{enumi})}
\item
  Gelfand 变换的核；
\item
  A 中满足 \(\mathrm{Sp}_{\mathrm{A}}^{\prime}(x)=\{0\}\) 的元素 x 所成的集合；
\item
  A 的拟幂零元所成的集合；
\item
  A 的根。
\end{enumerate}

我们分别用 \(N_{1}\)、\(N_{2}\)、\(N_{3}\)、\(N_{4}\) 表示这些集合。我们有 \(N_{1} = N_{2}\)（命题 6, a)），以及 \(N_{2} = N_{3}\) (I, p.~28, 推论 5)。按定义，集合 \(N_{4}\) 是 A 的正则极大理想的交；因而它是 A 的诸特征标的核的交 (I, p.~30, 定理 2)，后者等于 \(N_{1}\)。

注. --- 1) 一般说来，Gelfand 变换的像既不在 \(\mathcal{C}_{0}(\mathsf{X}(\mathsf{A}))\) 中闭，也不在 \(\mathcal{C}_{0}(\mathsf{X}(\mathsf{A}))\) 中稠密 (I, p.~193, 习题 7)。

\begin{enumerate}
\def\labelenumi{\arabic{enumi})}
\setcounter{enumi}{1}
\item
  Gelfand 变换的像分离 \(\mathsf{X}(\mathsf{A})\) 的点，因为若 \(\chi_1 \neq \chi_2\) 是 \(\mathsf{X}(\mathsf{A})\) 的元素，则存在 \(x \in \mathsf{A}\) 使得 \(\chi_1(x) \neq \chi_2(x)\)。
\item
  若 \(\chi \in \mathsf{X}(\mathsf{A})\)，则存在 Gelfand 变换的像中的一个元素，它在 \(\chi\) 处不取值 0。
\item
  若 A 有单位元，则 Gelfand 变换的像是 \(\mathsf{X}(\mathsf{A})\) 上连续函数代数的一个全子代数（命题 6, b)）。
\end{enumerate}

引理 2. --- 设 A 是一个交换 Banach 代数。设 M 是 X(A) 的一个子集。则 M 关于 Jacobson 拓扑是闭的，当且仅当对每个特征标 \(\chi \in \mathsf{X}(\mathsf{A}) = \mathsf{M}\)，存在 A 的一个元素 \(x\)，使得 \(\mathcal{G}(x)\) 在 M 上为零而在 \(\chi\) 处非零。

设 \(\Upsilon(\mathrm{M})\) 是 M 的诸元素的核的交。集合 M 关于 Jacobson 拓扑是闭的，当且仅当 \(\mathrm{M} = \mathrm{V}(\Upsilon(\mathrm{M}))\)（参见 I, p.~13 及 I, p.~30）。这个条件等价于说：M 的元素 \(\chi\) 恰好是在 \(\Upsilon(\mathrm{M})\) 上为零的那些特征标。因此，M 是闭的当且仅当对每个特征标 \(\chi \notin \mathrm{M}\)，存在 \(x \in \Upsilon(\mathrm{M})\) 使得 \(\chi(x) \neq 0\)。这可转写为 \(\mathcal{G}(x)(\chi) \neq 0\) 且 \(\mathcal{G}(x)|_{\mathrm{M}} = 0\)。

